\uuid{DhkJ}
\titre{Dérivée seconde ou plus}
\auteur{bodin ; hanani ; mzari}
\organisation{exo7}
\licence{CC BY-NC-SA 4.0}
\datecreate{2020-01-20}
\niveau{L1}
\module{Analyse}
\chapitre{Dérivabilité des fonctions réelles}
\sousChapitre{Autre}
\difficulte{1}
\theme{Dérivées, dérivée seconde ou plus}
\exo7qcmid{lille-1.derivees.010}

\contenu{
\begin{question}
\qclassification{124.00}
\qtags{motcle=dérivée seconde ou plus}
Soit $\displaystyle f(x)=\frac{1}{1+x}$. Quelles sont les bonnes réponses ?
\begin{answers}
    \good{$\displaystyle f''(x)=\frac{2}{(1+x)^3}$}
    \bad{$\displaystyle f''(x)=\frac{-2}{(1+x)^3}$}
    \bad{pour $n\in \Nn^*$, $\displaystyle f^{(n)}(x)=\frac{n}{(1+x)^{n+1}}$}
    \good{pour $n\in \Nn^*$, $\displaystyle f^{(n)}(x)=\frac{(-1)^nn!}{(1+x)^{n+1}}$}
\end{answers}
\begin{explanations}
On a $\displaystyle f'(x)=\frac{-1}{(1+x)^2}$, $\displaystyle f''(x)=\frac{2}{(1+x)^3}$ et l'on vérifie, par récurrence, que
$$\forall n\in \Nn^*,\; f^{(n)}(x)=\frac{(-1)^nn!}{(1+x)^{n+1}}.$$
\end{explanations}
\end{question}
}
